\documentclass[10pt]{article}

\usepackage[utf8]{inputenc}

\usepackage[T1]{fontenc}

\usepackage{amsmath}

\usepackage{amsfonts}

\usepackage{amssymb}

\usepackage[version=4]{mhchem}

\usepackage{stmaryrd}

\usepackage{graphicx}

\usepackage[export]{adjustbox}

\graphicspath{ {./images/} }

\usepackage{mathrsfs}



\begin{document}

характеристики к-рых не меняются с течением времени, т. е. пнвариантны относительно преобразования $X(t) \rightarrow$ $\rightarrow X(t+a)$, где $a$ - произвольное фиксированное число. Многомерные распределения вероятностей общего стационарного С. П. $X(t)$ не могут бить просто описаны, но для многих задач, касающихся таких процессов, достаточно знать лишь значения первых двух моментов $\mathrm{E} X(t)=m$ и $\mathrm{E} X(t) X(t+s)=B(s) \quad$ (так что здесь оказывается нужным лишь предположение о стационарности в широком смысле, т. е. о том, что от $t$ не зависят моменты $\mathrm{E} X(t)$ и $\mathrm{E} X(t) X(t+s))$. Существенно, что любой стационарный (или хотя бы стационарный в широком смысле) С. п. допускает спектральное разложение вида



\[

X(t)=\int_{-\infty}^{\infty} e^{i t \lambda} d Z(\lambda)

\]



где $Z(\lambda)$ - случайный процесс с некоррелированными шриращениями; отсюда, в частности, следует, что



\[

B(s)=\int_{-\infty}^{\infty} e^{i t \lambda} d F(\lambda)

\]



где $F(\lambda)$ - монотонно неубывающая спектральная функция С. п. $X(t)$. Спектральные разложения (5) и (6) лежат в основе решения задач о наилучпей (в смысле минимума среднеквадратитескоі̆ опибки) линейной экстраполяции, интерполяции и фильтрацип стационарных слутайных процессов.



Математич. теория С. п. включает также больпое число результатов, относящихся к ряду подклассов или, наоборот, обобщений перечисленных выше классов С. П. (в частности, к Маркова цеплм, диффузионныл процессам, ветьящимся процессам, мартингалам, случайным процессам со стационарными приращениями нек-рого порядка и др.).



Лит.: [1] С л у п к и й Е. Е., Избр. труды, М., 1960, с. $269-80 ;$ [2] Д у 6 Дж., Вероятностные процессы, пер. с Введение в теорию случайных процессов, 2 изд., М., 1977; [4] [5] K p a m e p $\Gamma$., П п д б е т т e p M., Стационарные случайные процессы, пер. с англ., М., 1969; [6] В е н т ц е л ь А. Д., Курс теории случайных процессов, М., 1975; [7] Р о з а н о в Ю. А., Случайные процессы, 2 изд., М., 1979; [8] И т о К., [9] С к о р о х о д А. В., Случайные процсссы с независимыми приращениями, М., 1964; [10] Д ы н к и н Е. Б., Марковские процессы, М., 1963, [11] И б р а г и м о в И. А., Р о з ан о в Ю.'А., Гауссовские случайные процессы, М., 1970; [12] 1963 . СЛУЧАЙЫИ ПРОЦЕСС ДИФФЕРЕНЦИРУЕМЫЙ - слутайный процесс $X(t)$ такой, что существует предел



\[

\lim _{\Delta t \rightarrow 0} \frac{X(t+\Delta t)-X(t)}{\Delta t}=X^{\prime}(t)

\]



называемый п р о и з в о д н о й с л у ч а й н о г 0 п р о ц е с с а $X(t)$; в зависимости от того, в каком смысле понимается этот предел, различают д и фф ер ен ц и ров ани е с ве роятностью 1 и



\includegraphics[max width=\textwidth, center]{2023_07_09_c24ce1a57a159ad2ec8fg-1(1)}

р а т и ч н о м. Условия дифференцируемости в среднем квадратичном естественно выражаются в терминах корреляционной функции



\[

B\left(t_{1}, t_{2}\right)=\mathrm{E} X\left(t_{1}\right) X\left(t_{2}\right)

\]



а именно $X^{\prime}(t)$ существует тогда и только тогда, когда существует предел



\[

B^{\prime \prime}\left(t_{1}, t_{2}\right)=

\]



$=\lim _{\substack{\Delta t_{1} \rightarrow 0 \\ \Delta t_{2} \rightarrow 0}} \frac{B\left(t_{1}+\Delta t_{1}, t_{2}+\Delta t_{2}\right)-B\left(t_{1}-\Delta t_{1}, t_{2}\right)-B\left(t_{1}, t_{2}+\Delta t_{2}\right)+B\left(t_{1}, t_{2}\right)}{\Delta t_{1} \Delta t_{2}}$.



Случайный процесс, шмеющий среднеквадратичную производную, является абсолютно непрерывным, точнее, шри каждом $t$ с вероятностью 1



\[

X(t)=X\left(t_{0}\right)+\int_{t_{0}}^{t_{1}} X^{\prime}(s) d s, \quad t \geqslant t_{0} .

\]



Достаточным условием того, чтобы существовал әквивалентный данному процесс с непрерывно диффееренцируемыми траекториями, может служить условие непрерывности его среднеквадратичной производной $X^{\prime}(t)$, имеюцей своей корреляционной функцией $B^{\prime \prime}\left(t_{1}, t_{2}\right)$. Для гауссовских процессов это условие является также необходимым.



\includegraphics[max width=\textwidth, center]{2023_07_09_c24ce1a57a159ad2ec8fg-1(2)}

в теорию случайных процессов, 2 изд., М., 1977. FО.' А. Розанов. чайный процесс с достаточно «простой» структурой, построенный по исходному процессу и содержащий всю требуемую инфформацию об этом процессе. С. п. о. использовались в задаче линейного прогноза стационарных случаӥных последовательностей, в нелинейных задачах статистики слутайных процессов и т. Д. (см. [1]-[3])



Понятие С. П. о. по-разному вводится в линейной и нелинейной теориях случайных процессов. В линейной теории (см. [4]) векторный слутайный процесс $x_{t}$ наз. о б н о в л ю ш и м п р о ц е с с о м для случайного процесса $\xi_{t}$ с E $\left|\xi_{t}\right|^{2}<\infty$, если $x_{t}$ имеет некоррелированные компоненты с некоррелированными приращениямп I



\[

H_{t}(\xi)=H_{t}(x) \text { для всех } t,

\]



где $H_{t}(\xi), H_{t}(x)$ - замкнутые в среднем квадратичном линейные оболочки всех значений $\xi_{s}, x_{s}, s \leqslant t$. Число компонент $N, N \leqslant \infty$, процесса $x_{t}$ наз. к р а т н о с т ь ю о б н о в л ю щ е г о п р о ц е с с а и определяется однозначно по процессу $\xi_{t}$. В случае дискретного времени $N=1$, а в случае непрерывного времени $N<\infty$ только при нек-рых специальных предположениях относительно корреляционной функции процесса $\xi_{t}$ (см. [4], [5]). В приложениях используется возможность представления $\xi_{t}$ в виде линейной комбинации знатений $x_{s}, s \leqslant t$.



В нелинейной теории (см. [5], [6]) обычно о 6 н о вля юши м случ а йным проц ес сом наз. винеровский процесс $x_{t}$ такой, что



\[

\mathscr{T}^{\xi}=\mathscr{F}^{x}

\]



где $\mathscr{F}^{\xi}, \mathscr{F}^{x}$ суть б-алгебры событий, порожденные значениями $\xi_{s}, x_{s}, s \leqslant t$. В случае когда $\xi_{t}, 0 \leqslant t \leqslant T$, является Ито процессом со стохастич. дифференциалом



\[

d \xi_{t}=\boldsymbol{a}(t) d t+d w_{t},

\]



винеровский процесс $\overline{w_{t}}$, определяемый равенством



\[

\bar{w}_{t}=\xi_{t}-\int_{0}^{t} \mathrm{E}\left(a(s) \mid \mathcal{F}_{s}^{\xi}\right) d s,

\]



является С. п. о. для $\xi_{t}$, если, напр.,



\[

\mathrm{E} \int_{0}^{T} a^{2}(s) d s<\infty

\]



и процессы $a$ и $w$ образуют гауссовскую систему (см. [6]).



Лum.: [1] К о л м о г о р о в А. Н., «Изв. АН СССР. Сep.



\includegraphics[max width=\textwidth, center]{2023_07_09_c24ce1a57a159ad2ec8fg-1}

"П а i l a t h T., "IEEE Ttans. "Inform. Theory", 1977, v. 1T-20 №2 2, р. 146 81; [4] Р о 3 а н о в Ю. А., Теория обновляющих процессов, М., 1974; [5] III и р я е в А. Н., в кн.: Тр. Школыч. 2 , Вильнюс, 1975 , с. $235-67 ;$ [6] Л и п ч е р Р. III., ІІІ ир я е в А. Н., Статистика случайных процессов, М., 1974.



СЛУЧАЙНЫЙ ПРОЦЕСС ОБОБІЕННЫЙ - сЛУчайный процесс $X$, зависящий от вепрерывного врецендо́го аргумента $t$ и такой, что его значения в фик-





\end{document}